\section{Base-Rate Study on Public GitHub Workflows}
\label{sec:realworld}

To estimate defect base rates on independently authored workflows---and measure
extractor fidelity on code we did not write---we mined public GitHub across all
four frameworks. The population is \emph{public GitHub files matching
framework-specific code searches}, not production deployments; the unit is one
source file mapped to one \emph{extracted file-level graph}, so a fully
executable-workflow-root unit is future work. All prevalence figures are
\emph{LLM-assisted exploratory estimates}, not ground truth.

\paragraph{Mining and corpus census.}
We mined GitHub on 2026-07-11 with eleven \texttt{gh search code} queries
(three per framework, two for CrewAI), retaining the top \textbf{40} results per
query in GitHub's default relevance order and cloning the first \textbf{180}
unique repositories in that order; within a cloned repository every
\texttt{.py} file was extracted (no per-file cap). After removing one repository
of our own test fixtures (8 files; a contamination detected during revision and
excluded from every number), the corpus is \textbf{922 extracted file-level
graphs} from 252 distinct repositories (LangGraph~400, AutoGen~359, CrewAI~113,
ADK~50), with provenance recorded as repository, commit SHA, and file path.
This is a \emph{search-ranked convenience corpus}, not a probability sample: the
ranking function is proprietary and non-stationary, the index denominator is
unpublished, and both cuts are deterministic in that ranking, so inclusion
probabilities are \emph{undefined} rather than merely unestimated. Every rate is
a property of this corpus and the instrument applied to it; we make no claim of
the form ``$X\%$ of agent workflows on GitHub.'' Mined graphs are small (median 2
non-sentinel nodes, mean~3.9) and predominantly static linear pipelines, and the
corpus is \emph{not} mostly applications: 32\% of files carry tutorial/test/demo
names, 170 graphs are topological duplicates, and a source-level classification
of the validation sample found 44\% tutorial, 17\% toy demo, 7\% test, and 33\%
application-like code.

\paragraph{Reference graphs and validation.}
On a 119-graph validation sample from 87 repositories (LangGraph~40, CrewAI~20,
AutoGen~35, ADK~24) we ran two independent LLM-agent passes (Claude Opus~4.8
under a multi-agent harness; all outputs released). First, an agent
reconstructed a reference graph from source alone, blind to the extractor's
output and following its node-id conventions; we score node/edge
precision--recall and node-kind accuracy against these \emph{LLM-reconstructed
reference graphs} (confidence: 77 high, 33 medium, 9 low). Second, every flag was
triaged from source
(\textsc{real}/\textsc{extraction-artifact}/\textsc{intentional}) and an
adversarial verifier attempted to overturn each label (4.6\% overturned;
unresolved labels recorded as \textsc{arguable}). Reference graphs are LLM
reconstructions, not human annotations, and both passes share a model family.
Results are stable on the high-confidence subset, every triage label carries a
source-line justification that survived an adversarial pass, and a re-audit
overturned one label of 16 re-examined.
\textbf{Two-annotator independent human validation was not performed; it is
required future work, and the protocol and samples are staged.} We report Wilson
95\% CIs and repository-clustered bootstrap CIs (10{,}000 resamples), since
workflows from one repository are not independent.

\paragraph{How the validation sample was selected.}
The 119 graphs were chosen under fixed per-framework quotas \emph{without
randomization}: the identifiers exist only as literal lists in the released
harness scripts, with no seed or reproducible selection rule. We therefore do not
claim the sample is random within framework, nor call it stratified.
Consequently the framework-share reweighting is a
\emph{corpus-share-weighted descriptive statistic}---the rate the corpus would
show if each framework's sampled rate held across its members, correcting the
sample's ${\sim}5\times$ over-representation of ADK---not a design-based
estimate; and the intervals quantify sampling noise \emph{within this corpus}
only. We apply no finite-population correction, since an FPC presupposes the
simple random sampling this design does not provide.

\begin{table}[t]
\centering\small
\caption{Extractor fidelity on the \emph{matched, provenance-disambiguated} set
($n{=}106$): both instruments scored on identical graphs, colliding records
excluded. Edge P and kind are identical between instruments except where
shown---the corrected instrument changes \emph{exactly one} framework.}
\label{tab:realworld_fidelity}
\resizebox{\columnwidth}{!}{%
\begin{tabular}{@{}lrrrrr@{}}
\toprule
\textbf{Framework} & \textbf{$n$} & \textbf{Edge P} & \textbf{Edge R$_{\mathrm{v1}}$} & \textbf{Edge R$_{\mathrm{v2}}$} & \textbf{Kind} \\
\midrule
LangGraph & 38 & 0.954 & 0.678 & \textbf{0.836} & 0.659/0.670 \\
CrewAI    & 20 & 0.700 & 0.692 & 0.692 & 0.929 \\
AutoGen   & 30 & 0.966 & 0.778 & 0.778 & 0.890 \\
ADK       & 18 & 0.394 & 0.343 & 0.343 & 0.920 \\
\midrule
\textbf{Overall} & \textbf{106} & \textbf{0.814} & \textbf{0.652} & \textbf{0.709} & \textbf{0.820/0.824} \\
\bottomrule
\end{tabular}}
\end{table}

\paragraph{A keying defect.}
We report an instrumentation error found while preparing this comparison,
because it changes the numbers. Graphs were keyed by
repository plus file \emph{basename}, which is not unique: \textbf{76} slugs
collide, and for those the two instruments scored \emph{different source files}
against the same reference graph. The v1 miner resolved collisions last-wins
without recording which file survived, so no colliding pair is repairable (all
58 affected v1 records; supplement) and every colliding graph is dropped:
\textbf{9} inside the fidelity set, cutting it from 115 matched graphs to the
$n{=}106$ of Table~\ref{tab:realworld_fidelity}, with 4 more of the 119 sampled
unavailable at re-mining. The cost is not uniform---ADK loses 5 of 23 versus
0--6\% elsewhere---so ADK's row is the least stable. The keying is fixed in the
released instrument.

\paragraph{Extraction fidelity is the bottleneck---for structural checks.}
Node detection stays strong and stable across instruments (v2 overall
$P{=}0.914$, $R{=}0.868$; v1 $0.914/0.862$), but edge recall and node-kind
accuracy are the bottleneck and are strongly framework-dependent
(Table~\ref{tab:realworld_fidelity}): the corrected instrument recovers only
0.34 of ADK's edges---whose nested \texttt{SequentialAgent}/
\texttt{ParallelAgent} composition, and AutoGen's participant lists with dynamic
speaker selection, encode topology \emph{implicitly}---versus 0.84 on LangGraph's
explicit edges. On LangGraph the dominant remaining error is instead
node-kind (\texttt{tool}$\rightarrow$\texttt{llm}: nodes invoking tools in their
body with no declared binding; kind accuracy 0.67), the direct cause of the
risk-aware gate's blindness below. This is a property of the static AST fallback
used at scale, not of the runtime extractors reading the compiled framework
object (for ADK the runtime graph is far richer, 16 vs.\ ${\sim}5$ edges);
realizing that ceiling would require sandboxed execution of untrusted code, which
we do not attempt.

A corrected-instrument ablation demonstrates the diagnosis \emph{causally} for
the structural checks. The correction is not a single fix: the instrument
bundles \emph{three} behavioral changes---a LangGraph three-argument
\texttt{add\_conditional\_edges} fix, a LangGraph declared-tool metadata fix,
and ADK \texttt{LoopAgent} support. Re-extracting at identical
pinned commits with the check algorithms untouched, only \emph{one} has a
measurable effect: LangGraph edge recall rises $0.678\rightarrow0.836$, and of
the 54 true edges newly recovered, \textbf{54/54} are targets of a three-argument
\texttt{add\_conditional\_edges} path map across 24 call sites. ADK, AutoGen, and
CrewAI are bit-identical between instruments
(Table~\ref{tab:realworld_fidelity}), so the \texttt{LoopAgent} fix cannot be
credited with an improvement. On the matched collision-free corpus ($n{=}853$)
the correction removes roughly a quarter of the structural flags
($299\rightarrow229$ flagged workflows; unreachable-exit flags
$81\rightarrow29$), confirming extraction as the binding constraint \emph{for
the structural checks}---a scope the attribution below makes precise, and which
does not extend to the policy checks. Both instruments are released.

\paragraph{Flag precision on extracted graphs.}
Distinct from prevalence, what fraction of the checker's \emph{flags} is
genuine? The as-mined study raised 186 flags on 114 file-level graphs;
adversarial triage confirmed 2 as genuine (both missing human gates),
classifying 42\% as extraction artifacts and 50\% as intentional design.
That is a workflow-level PPV of $2/114 = 1.75\%$ $[0.48,6.17]\%$ and a flag-level
PPV of $2/186 = 1.08\%$ $[0.30,3.84]\%$; deduping correlated flags to root cause
($186\rightarrow161$) gives $1.24\%$, and counting all 12 \textsc{arguable}
labels pessimistically bounds the rate at $7.53\%$, so identification lies in
$[1.08\%,7.53\%]$. \emph{None of the 72 structural flags was confirmed genuine}
(CI $[0,5.1]\%$), and structural flags are almost entirely a LangGraph
phenomenon---partly \emph{by construction}, since the other three extractors
emit connected topologies and cannot produce a dead end or unreachable exit.
This as-mined universe is \emph{distinct} from the reconstructed-graph re-run
below; the two are never pooled.

\paragraph{Two failure modes.}
A precision number says findings are non-actionable but not \emph{why}, and the
answer decides what to fix. We attribute each of the 186 flags to one of four
causes: \textsc{extractor} (the graph misrepresents the code); \textsc{checker}
(check logic wrong given a correct graph); \textsc{policy-spec} (graph and check
both correct, but the policy does not apply---a blunt human-presence requirement
on a read-only tutorial); and \textsc{label} (disputed judgment). One caveat: our
triage vocabulary asked ``is this a real defect?'', not ``whose fault is this
flag?'', so \textsc{checker} is not separable from \textsc{policy-spec} and is
folded into it---nonzero but unquantified; both are check-side, so the comparison
is unaffected. The result does not support one cause: extraction artifacts
account for $42.5\%$ of flags $[35.6,49.7]$, policy misspecification for $50.0\%$
$[42.9,57.1]$---overlapping intervals in which extraction is, if anything, the
\emph{smaller} term.

\paragraph{The decisive test.}
Attribution rests on labels, so we ran the symmetric ablation: hold the checks
fixed, vary only the graph, using the 119 reference graphs as a near-oracle
extractor. Flags fall from 184 to 122---an oracle removes \textbf{62 flags,
roughly a third}, leaving \textbf{${\ge}120$ non-actionable findings standing} on
119 workflows, a check-side floor ${\approx}1.9\times$ the
extraction-attributable share. The other knob is far more effective: holding the
graph fixed and swapping the blunt policy for the risk-aware gate moves firings
from \textbf{113 to 0}. The causes separate by check family: among
non-actionable \emph{structural} flags ($n{=}72$), $98.6\%$ $[92.5,99.8]$ are
extraction-attributable and $0\%$ $[0,5.1]$ check-side, whereas among
non-actionable \emph{human-gate} flags ($n{=}112$) only $7.1\%$ $[3.7,13.5]$ are
extraction-attributable against $83.0\%$ $[75.0,88.9]$ check-side---$92\%$ of
attributed flags, the rest disputed labels. Because
human-gate checks are $61\%$ of all flags, the aggregate goes against a pure
extraction story even though extraction dominates the structural checks. Better
extractors fix reachability checks; only per-workflow policy instantiation fixes
safety checks; neither alone suffices.

% Attribution table moved to the supplement (page limit); all figures stated inline.

\paragraph{Defining a defect.}
Flag triage cannot see false negatives: a defect the extractor never surfaces is
absent from the flag universe entirely. We therefore re-ran every check on the
119 reference graphs---which can express dead ends and unreachable exits for all
four frameworks---and adversarially triaged the results, recovering two genuine
defects the extracted study missed. A \emph{structural defect} is a topology
error; a \emph{human-gate policy violation} is a normative judgment that a
side-effecting tool is reachable with no intervening \textsc{human} checkpoint,
which a permissive policy would not yield. We report each as a \textbf{distinct
estimand} and their union only as a heterogeneous secondary summary.

\paragraph{Three estimands and reweighting.}
Table~\ref{tab:estimands} reports all three under the uniform policy. The one
structural defect is a router mis-wiring in an ADK SQL agent, $1/119$
$[0.15,4.61]\%$; the human-gate estimand is the re-audited $12/119$
$[5.86,16.80]\%$; and the composite unions them---the structural case is
\textsc{arguable}, not a violation, under HGP-1, so the union is a strict
addition---giving $13/119$ $[6.50,17.80]\%$. Per-framework violations run
LangGraph $6/40$, AutoGen $5/35$, CrewAI $0/20$, ADK $1/24$. Reweighting each
stratum to its corpus share (clustered bootstrap, no finite-population
correction) \emph{raises} rather than lowers these estimates, because LangGraph
and AutoGen carry both the highest violation rates and $82\%$ of corpus weight:
human-gate $12.30\%$ $[4.70,21.34]$, composite $12.52\%$ $[4.93,21.56]$. The
weakest link is CrewAI's $0\%$, resting on a single audited workflow.

The revision is \emph{not} additive. Of the three cases the two-check procedure
reported, HGP-1 retains only one; the other two become \textsc{arguable}, their
side-effecting callees living in modules the single-file mining unit never
captured. The uniform policy thus retains 1 of 3 and finds 11 new violations. We
do not re-estimate the earlier application-versus-tutorial stratum association,
which described a superseded four-case universe.

\begin{table}[t]
\centering\small
\caption{Three defect estimands on the 119-graph validation sample (LLM-assisted
exploratory estimates), with Wilson 95\% CIs (\%) and corpus-share-weighted
estimates (clustered bootstrap, no FPC). The human-gate row is the uniform HGP-1
re-audit, superseding the $3/119$ a two-check procedure reported.}
\label{tab:estimands}
\resizebox{\columnwidth}{!}{%
\begin{tabular}{@{}llll@{}}
\toprule
\textbf{Estimand} & \textbf{Sample} & \textbf{Wilson 95\%} & \textbf{Weighted} \\
\midrule
Structural defect      & 1/119 (0.84\%)  & [0.15,\,4.61]  & 0.23\% [0.00,\,0.80] \\
Human-gate (HGP-1)     & 12/119 (10.1\%) & [5.86,\,16.80] & 12.30\% [4.70,\,21.34] \\
\quad + \textsc{arguable} & 23/119 (19.3\%) & [13.24,\,27.34] & 17.74\% [9.52,\,27.01] \\
Composite$^{\dagger}$  & 13/119 (10.9\%) & [6.50,\,17.80] & 12.52\% [4.93,\,21.56] \\
\bottomrule
\end{tabular}}\\[2pt]
{\footnotesize $^{\dagger}$Heterogeneous secondary summary (union of the two
primary estimands).}
\end{table}

\paragraph{One policy, applied uniformly.}
The three policy cases above came from \emph{two} different checks (blunt
\texttt{human\_presence} and path-sensitive \texttt{human\_gate\_coverage}) over
two graph populations. That is not one estimand, and it cannot see a violation
neither check flags. We therefore pre-registered a single operational policy
(HGP-1; in the supplement) and applied it uniformly to every sampled workflow
with a side-effecting tool at \emph{source} level---all 32---regardless of flag
universe. HGP-1 fixes the side-effecting categories by what the callee does in
source rather than by name; requires a gate to be an explicit checkpoint or a
blocking confirmation whose negative branch exists and skips the effect, and to
\emph{dominate} the call site; and records undecidable cases as
\textsc{arguable}, never silently as compliant. Demo, tutorial, and test status
are explicitly \emph{not} exemptions; genuine mocks and visibly sandboxed effects
are.

The audit returns \textbf{12 violations, 9 compliant, 11 arguable}: $12/119 =
10.1\%$ $[5.9,16.8]$ against the reported $3/119$---a fourfold revision whose
interval excludes the old estimate; \textsc{arguable} counted pessimistically
bounds it at $23/119$. The uncertainty points upward: $12/119$ is a \emph{lower
bound}, since several \textsc{arguable} cases (a payment pipeline whose prompts
forbid confirmation four times over) would likely resolve toward violation given
repository-level rather than single-file source. The 87 workflows without
side-effecting tools were screened by the same source-level pass, and a re-sweep
for effect signatures surfaced no missed violation. Three of the 12 lie
\emph{outside both flag universes}---no triage entry, neither check
fired---confirming that false negatives exist beyond the union of the flag sets,
where a flag-triage study can never find them.

\paragraph{Placebo gates: the reference graphs overstate oversight.}
The most consequential finding is not a count. All seven side-effecting
workflows that \emph{both} checks cleared were cleared because the reference
graph held a node typed \textsc{human}---and on source inspection \emph{none is
a genuine dominating gate}: five are placebo or inert, one does not dominate the
call site, one is unverifiable. One exemplar declares six
\texttt{*\_human\_approval} nodes whose body reads a constructor flag hard-set
to \texttt{True}: a checkpoint that cannot block, typed \textsc{human} by both
extractor and reconstruction because it is named and shaped like one. Reference
graphs assign \textsc{human} from naming and framework type, not from whether the
body can block, so \emph{every} human-gate estimand computed over
them---including the risk-aware one---is systematically biased \textbf{downward}.
The fidelity metric compounds it: for two of these workflows it records
human-node recall $1.0$, crediting the extractor with recovering oversight that
does not exist. A near-oracle graph is not a semantic oracle, and node
\emph{kinds} inherit whatever the source calls its nodes.

\paragraph{Risk-aware gating.}
The blunt human-presence check fires on 91\% of mined graphs with only 2 of 113
firings genuine---uninformative as a detector, motivating the risk-aware
\emph{human-gate coverage} check: flag only when a sensitive tool
(destructive/write/exec/communication/financial, by a keyword lexicon over bound
tool names) is reachable without a \textsc{human} gate. On the 922 mined graphs
it fires zero times---a statement about extracted \emph{declarations}, not the
code: no mined graph declares a sensitive binding, while source review finds
side-effecting tools in 32 of 119 sampled workflows, invisible to the extractor
through the \texttt{tool}$\rightarrow$\texttt{llm} kind error. On reconstructed
graphs it fires three times. The lexicon is \emph{correct relative to its
declared-tool abstraction}; its recall is bounded by extraction fidelity, its
precision---per the audit above---by whether a \textsc{human} node can block.

\paragraph{Threats to validity.}
\label{par:threats}
\emph{Construct:} the composite unions a structural property with a policy
judgment, so we report the two separately. \emph{Internal:} reference graphs and
triage labels come from LLM passes sharing a model family, so correlated errors
are possible; confidence stratification, adversarial overturn rates, and a
re-audit mitigate but do not replace independent human judgment.
\emph{External:} the population is public GitHub files and the unit a file-level
graph, so results may not transfer to executable roots or private code.
\emph{Sampling:} the frame is a convenience corpus with undefined inclusion
probabilities and the sample was quota-selected without randomization, so
reweighting yields a descriptive statistic, not a design-based estimate.
\emph{Label reliability:} no human validation has been performed; a two-annotator
protocol with staged samples is future work.
\emph{Reference-graph semantics:} node kinds are assigned from naming, not from
whether a body has the effect its kind implies---the placebo-gate audit shows
this biases human-gate estimands downward. Every prevalence number is accordingly
an LLM-assisted exploratory estimate, not ground truth.
